% !TEX root = 第一章_多项式(答案版).tex
\setcounter{chapter}{0}
\pagestyle{fancy}

\chapter{多项式}

\section{整除}

%例题1.1.1
\begin{example}
设
\[
  f(x)=x^6-1,\qquad g(x)=x^3+x+1,
\]
求带余除法中满足
\[
  f(x)=g(x)q(x)+r(x),\qquad \deg r<\deg g
\]
的 $q(x),r(x)$.
\end{example}
\exampleblank

\begin{solution}
直接作带余除法:
\[
  x^6-1=(x^3+x+1)(x^3-x-1)+(x^2+2x).
\]
因此
\[
  q(x)=x^3-x-1,\qquad r(x)=x^2+2x.
\]
\end{solution}

%例题1.1.2
\begin{example}
求 $a,b$,使得
\[
  x^2-2ax+2\mid x^4+3x^2+ax+b.
\]
\end{example}
\exampleblank

\begin{solution}
模 $x^2-2ax+2$ 有
\[
  x^2\equiv 2ax-2.
\]
逐次降幂可得
\[
  x^4+3x^2+ax+b
  \equiv (8a^3-a)x+b-8a^2-2.
\]
故整除的充要条件为
\[
  8a^3-a=0,\qquad b-8a^2-2=0.
\]
于是
\[
  (a,b)=(0,2),\quad
  \left(\dfrac{\sqrt2}{4},3\right),\quad
  \left(-\dfrac{\sqrt2}{4},3\right).
\]
\end{solution}

%例题1.1.3
\begin{example}
设
\[
  f(x)=2x^4-3x^3+4x^2+ax+b,
  \qquad
  g(x)=x^2-3x+1.
\]
若 $f(x)$ 除以 $g(x)$ 后余式为 $25x-5$,求 $a,b$.
\end{example}
\exampleblank

\begin{solution}
由 $x^2\equiv3x-1\pmod{g(x)}$ 逐次降幂,得到
\[
  f(x)\equiv (a+30)x+b-11\pmod{g(x)}.
\]
与余式 $25x-5$ 比较系数,得
\[
  a=-5,\qquad b=6.
\]
\end{solution}

%例题1.1.4
\begin{example}[\markstar]
利用综合除法,将
\[
  2x^4+x^3+7x^2-8x+14
\]
表示成
\[
  ax(x-1)(x-2)(x-3)
  +bx(x-1)(x-2)
  +cx(x-1)+dx+e,
\]
求 $a,b,c,d,e$.
\end{example}
\exampleblank

\begin{solution}
这是以节点 $0,1,2,3$ 为基点的 Newton 型展开.依次作综合除法可得:
\[
\begin{array}{c|rrrrr}
0&2&1&7&-8&14\\
 & &0&0&0&0\\ \hline
 &2&1&7&-8&14
\end{array}
\]
故 $e=14$;再对商的相应差分依次在 $1,2,3$ 处作综合除法,可得
\[
  d=2,\qquad c=24,\qquad b=13,\qquad a=2.
\]
因此
\[
\begin{aligned}
2x^4+x^3+7x^2-8x+14
={}&2x(x-1)(x-2)(x-3)\\
&+13x(x-1)(x-2)+24x(x-1)+2x+14.
\end{aligned}
\]
\end{solution}

%例题1.1.5
\begin{example}[\markstar]
证明:
\[
  x^d-1\mid x^n-1
\]
的充要条件是 $d\mid n$.
\end{example}
\exampleblank

\begin{solution}
若 $d\mid n$,设 $n=qd$,则
\[
  x^n-1=(x^d)^q-1
  =(x^d-1)\left(1+x^d+\cdots+x^{(q-1)d}\right).
\]
故 $x^d-1\mid x^n-1$.

反之,设 $n=qd+r$,其中 $0\le r<d$.模 $x^d-1$ 有 $x^d\equiv1$,故
\[
  x^n-1\equiv x^r-1\pmod{x^d-1}.
\]
若 $x^d-1\mid x^n-1$,则 $x^d-1\mid x^r-1$.当 $0<r<d$ 时,
$\deg(x^r-1)<\deg(x^d-1)$,不可能被其整除.因此 $r=0$,即 $d\mid n$.
\end{solution}

%例题1.1.6
\begin{example}[\markstar\ 不可约多项式的素性]
设 $f(x)\in F[x]$ 是一个次数大于零的多项式.证明:$f(x)$ 不可约的充要条件是:
若
\[
  f(x)\mid g(x)h(x),
\]
则必有 $f(x)\mid g(x)$ 或 $f(x)\mid h(x)$.
\end{example}
\exampleblank

\begin{solution}
若 $f$ 不可约,而 $f\mid gh$.若 $f\nmid g$,则由 $f$ 不可约可知
\[
  (f,g)=1.
\]
由 Bézout 等式存在 $u,v\in F[x]$ 使
\[
  uf+vg=1.
\]
乘以 $h$ 得
\[
  ufh+vgh=h.
\]
左端两项均被 $f$ 整除,故 $f\mid h$.

反之,若题述素性成立而 $f$ 可约,则可写成
\[
  f=gh,
  \qquad
  0<\deg g,\deg h<\deg f.
\]
显然 $f\mid gh$,但 $f\nmid g$ 且 $f\nmid h$,矛盾.因此 $f$ 不可约.
\end{solution}


\newpage

\section{最大公因式}

\subsection{最大公因式的计算与性质}

%例题1.2.1
\begin{example}
设
\[
\begin{aligned}
  f(x)&=x^4+x^3-3x^2-4x-1,\\
  g(x)&=x^3+x^2-x-1,
\end{aligned}
\]
求 $u(x),v(x)$ 使得
\[
  u(x)f(x)+v(x)g(x)=(f(x),g(x)).
\]
\end{example}
\exampleblank

\begin{solution}
用扩展 Euclid 算法可得首一最大公因式
\[
  (f,g)=x+1,
\]
并且
\[
  \left(\dfrac13-\dfrac23x\right)f(x)
  +\left(\dfrac23x^2-\dfrac13x-\dfrac43\right)g(x)
  =x+1.
\]
因此可取
\[
  u(x)=\dfrac{1-2x}{3},
  \qquad
  v(x)=\dfrac{2x^2-x-4}{3}.
\]
\end{solution}

%例题1.2.2
\begin{example}
求下列多项式的最大公因式:
\begin{enumerate}
  \item $f(x)=x^4-x^3-4x^2+4x+1$,$g(x)=x^2-x-1$;
  \item $f(x)=x^4+3x^3-x^2-4x-3$,$g(x)=3x^3+10x^2+2x-3$.
\end{enumerate}
\end{example}
\exampleblank

\begin{solution}
由 Euclid 算法:
\begin{enumerate}
  \item $(f,g)=1$;
  \item $(f,g)=x+3$.
\end{enumerate}
\end{solution}

%例题1.2.3
\begin{example}[\markstar]
设 $f(x),g(x)\in F[x]$,$a,b,c,d\in F$,且
\[
  \begin{vmatrix}
    a&b\\
    c&d
  \end{vmatrix}\neq0.
\]
证明:
\[
  (f,g)=(af+bg,cf+dg).
\]
\end{example}
\exampleblank

\begin{solution}
记
\[
  F_1=af+bg,\qquad G_1=cf+dg.
\]
显然 $(f,g)$ 同时整除 $F_1,G_1$,故
\[
  (f,g)\mid(F_1,G_1).
\]
又因 $ad-bc\neq0$,有
\[
  f=\dfrac{dF_1-bG_1}{ad-bc},
  \qquad
  g=\dfrac{-cF_1+aG_1}{ad-bc}.
\]
于是 $(F_1,G_1)$ 同时整除 $f,g$,故
\[
  (F_1,G_1)\mid(f,g).
\]
取首一最大公因式即得二者相等.
\end{solution}

%例题1.2.4
\begin{example}
设
\[
  f(x)=x^3+(1+k)x^2+2x+2l,
  \qquad
  g(x)=x^3+kx^2+l,
\]
且 $(f,g)$ 是二次多项式,求 $k,l$.
\end{example}
\exampleblank

\begin{solution}
因为
\[
  f-g=x^2+2x+l,
\]
而 $(f,g)$ 为二次首一多项式,所以必有
\[
  (f,g)=x^2+2x+l.
\]
令其整除 $g$,作带余除法并令余式为零,得到
\[
  (k,l)=(2,0)\quad\text{或}\quad(3,-2).
\]
\end{solution}

%例题1.2.5
\begin{example}[\markstar]
证明:若 $(f,g)=1$ 且 $f\mid gh$,则 $f\mid h$.
\end{example}
\exampleblank

\begin{solution}
由 $(f,g)=1$,存在 $u,v\in F[x]$ 使
\[
  uf+vg=1.
\]
乘以 $h$:
\[
  ufh+vgh=h.
\]
由 $f\mid ufh$ 且 $f\mid vgh$,故 $f\mid h$.
\end{solution}

%例题1.2.6
\begin{example}[\markstar]
证明:若 $f_1\mid g$,$f_2\mid g$ 且 $(f_1,f_2)=1$,则
\[
  f_1f_2\mid g.
\]
\end{example}
\exampleblank

\begin{solution}
设 $g=f_1q$.由 $f_2\mid g=f_1q$ 且 $(f_1,f_2)=1$,由上一题得 $f_2\mid q$.
故 $q=f_2h$,于是
\[
  g=f_1f_2h.
\]
\end{solution}

%例题1.2.7
\begin{example}[\markstar]
证明:若 $(f,g)=d$,则
\[
  (tf,tg)=td.
\]
\end{example}
\exampleblank

\begin{solution}
写成
\[
  f=df_1,\qquad g=dg_1,
  \qquad (f_1,g_1)=1.
\]
于是
\[
  tf=(td)f_1,\qquad tg=(td)g_1.
\]
故 $td$ 是二者的公因式.又若 $c$ 是 $tf,tg$ 的任一公因式,利用
$(f_1,g_1)=1$ 的 Bézout 等式可得 $c\mid td$.因此最大公因式为 $td$(相差非零常数时取相应的首一化).
\end{solution}

%例题1.2.8
\begin{example}
设 $f,g$ 是次数不小于 $1$ 的互素多项式.证明:存在唯一的 $u,v$ 使
\[
  uf+vg=1,
\]
且
\[
  \deg u<\deg g,
  \qquad
  \deg v<\deg f.
\]
\end{example}
\exampleblank

\begin{solution}
由 $(f,g)=1$,Bézout 等式保证存在 $u_0,v_0$ 使
\[
  u_0f+v_0g=1.
\]
将 $u_0$ 除以 $g$:
\[
  u_0=qg+u,
  \qquad
  \deg u<\deg g.
\]
代入后令 $v=v_0+qf$,便有 $uf+vg=1$.由
\[
  vg=1-uf
\]
及次数比较可得 $\deg v<\deg f$.

若还有另一组 $u_1,v_1$ 满足相同次数限制,则
\[
  (u-u_1)f=-(v-v_1)g.
\]
由 $(f,g)=1$ 得 $g\mid u-u_1$.但 $\deg(u-u_1)<\deg g$,故 $u=u_1$,进而 $v=v_1$.
\end{solution}

\subsection{中国剩余定理(孙子定理)}

\begin{theorem}[\markstar\ 中国剩余定理]
设 $f_1(x),\ldots,f_n(x)$ 两两互素,$a_1(x),\ldots,a_n(x)$ 为任意多项式,则存在多项式 $g(x)$ 使
\[
  g(x)\equiv a_i(x)\pmod{f_i(x)},
  \qquad i=1,2,\ldots,n.
\]
\end{theorem}

\begin{proof}
对每个 $i$,令
\[
  F_i(x)=\prod_{j\ne i}f_j(x).
\]
由于 $f_i$ 与 $F_i$ 互素,存在 $u_i,v_i$ 使
\[
  u_iF_i+v_if_i=1.
\]
令 $e_i=u_iF_i$,则
\[
  e_i\equiv1\pmod{f_i},
  \qquad
  e_i\equiv0\pmod{f_j}\quad(j\ne i).
\]
于是
\[
  g=\sum\limits_{i=1}^n a_ie_i
\]
满足全部同余式.
\end{proof}

%例题1.2.9
\begin{example}
求次数最小的多项式 $f(x)$,使
\[
\begin{cases}
  f(x)\equiv6\pmod{x+1},\\
  f(x)\equiv3x\pmod{x^2+x+1},\\
  f(x)\equiv(x-1)^2\pmod{2x^3+1}.
\end{cases}
\]
\end{example}
\exampleblank

\begin{solution}
三个模多项式两两互素.设 $\deg f<1+2+3=6$,令
\[
  f(x)=c_0+c_1x+\cdots+c_5x^5.
\]
分别取余并比较系数,解得
\[
  c_0=1,\quad c_1=0,\quad c_2=1,\quad c_3=0,\quad c_4=4,\quad c_5=0.
\]
因此次数最小的解为
\[
  \boxed{f(x)=4x^4+x^2+1}.
\]
\end{solution}

%例题1.2.10
\begin{example}
找出所有满足
\[
\begin{cases}
  x\equiv2\pmod5,\\
  x\equiv1\pmod7,\\
  x\equiv5\pmod{11}
\end{cases}
\]
的整数 $x$.
\end{example}
\exampleblank

\begin{solution}
由中国剩余定理模 $5\cdot7\cdot11=385$ 有唯一解类.直接求得
\[
  x\equiv302\pmod{385}.
\]
故全部整数解为
\[
  \boxed{x=302+385k,\qquad k\in\mathbb{Z}}.
\]
\end{solution}

\subsection{不可约因式的素性}

%例题1.2.11
\begin{example}
设 $(f,g)=1$ 且 $(f,h)=1$,证明
\[
  (f,gh)=1.
\]
\end{example}
\exampleblank

\begin{solution}
若不可约多项式 $p$ 同时整除 $f$ 与 $gh$,由不可约因式的素性,$p\mid g$ 或 $p\mid h$,
从而与 $(f,g)=1$ 或 $(f,h)=1$ 矛盾.因此 $f$ 与 $gh$ 无公共不可约因式,即 $(f,gh)=1$.
\end{solution}

%例题1.2.12
\begin{example}
证明:对任意正整数 $n$,
\[
  (f,g)^n=(f^n,g^n).
\]
\end{example}
\exampleblank

\begin{solution}
写出标准分解
\[
  f=c\prod p_i^{\alpha_i},
  \qquad
  g=d\prod p_i^{\beta_i},
\]
其中对缺失因子允许相应指数为 $0$.则
\[
  (f,g)=\prod p_i^{\min(\alpha_i,\beta_i)}.
\]
故
\[
  (f^n,g^n)
  =\prod p_i^{\min(n\alpha_i,n\beta_i)}
  =\left(\prod p_i^{\min(\alpha_i,\beta_i)}\right)^n
  =(f,g)^n.
\]
\end{solution}

%例题1.2.13
\begin{example}
设 $(f,g)=1$,证明
\[
  (fg,f+g)=1.
\]
\end{example}
\exampleblank

\begin{solution}
若不可约多项式 $p\mid fg$ 且 $p\mid f+g$,则由 $p\mid fg$ 知 $p\mid f$ 或 $p\mid g$.
若 $p\mid f$,再由 $p\mid f+g$ 得 $p\mid g$;另一种情形同理.均与 $(f,g)=1$ 矛盾.
\end{solution}

%例题1.2.14
\begin{example}
设
\[
  (f_i,g_j)=1,
  \qquad
  i=1,\ldots,s,\quad j=1,\ldots,t.
\]
证明
\[
  (f_1\cdots f_s,g_1\cdots g_t)=1.
\]
\end{example}
\exampleblank

\begin{solution}
若某不可约多项式 $p$ 同时整除两边的两个乘积,则由素性可得某个 $i,j$ 使
\[
  p\mid f_i,
  \qquad
  p\mid g_j,
\]
与 $(f_i,g_j)=1$ 矛盾.
\end{solution}

\subsection{最大公因式的不变性}

\begin{theorem}[\markstar]
设 $f(x),g(x)\in K[x]$,且数域 $F\supseteq K$.则 $f,g$ 在 $K[x]$ 中的首一最大公因式,
等于它们在 $F[x]$ 中的首一最大公因式.特别地,两个多项式的互素性不随数域扩大而改变.
\end{theorem}

%例题1.2.15
\begin{example}
设 $f(x)\in\mathbb{Q}[x]$,且 $1+\sqrt3$ 为 $f(x)$ 的根.证明 $1-\sqrt3$ 也是 $f(x)$ 的根.
\end{example}
\exampleblank

\begin{solution}
$1+\sqrt3$ 在 $\mathbb{Q}$ 上的最小多项式为
\[
  (x-1)^2-3=x^2-2x-2,
\]
它在 $\mathbb{Q}[x]$ 中不可约.既然 $f(1+\sqrt3)=0$,则最小多项式整除 $f$,
故其另一根 $1-\sqrt3$ 也必为 $f$ 的根.
\end{solution}

%例题1.2.16
\begin{example}
设 $f(x)\in\mathbb{Q}[x]$,且
\[
  a\sqrt b+c\sqrt d
\]
为 $f(x)$ 的根,其中 $b,d$ 不是平方数.证明
\[
  a\sqrt b-c\sqrt d,
  \quad
  -a\sqrt b+c\sqrt d,
  \quad
  -a\sqrt b-c\sqrt d
\]
也为 $f(x)$ 的根.
\end{example}
\exampleblank

\begin{solution}
按题面当前条件,此结论并非总成立.若补充通常所需条件
\[
  [\mathbb{Q}(\sqrt b,\sqrt d):\mathbb{Q}]=4
\]
(例如 $b,d,bd$ 均不是有理数平方,且 $ac\ne0$),则
$\mathbb{Q}(\sqrt b,\sqrt d)$ 的四个 $\mathbb{Q}$-自同构分别独立改变
$\sqrt b,\sqrt d$ 的符号.将这些自同构作用于
\[
  f(a\sqrt b+c\sqrt d)=0
\]
即可得到另外三个数也都是 $f$ 的根.

原题条件不足,具体反例与说明见项目根目录的《解答问题记录.md》.
\end{solution}

%例题1.2.17
\begin{example}
设 $f(x)\in\mathbb{Q}[x]$,$d\in\mathbb{Q}$,且 $\sqrt[3]{d}$ 为无理数并且是 $f(x)$ 的一个根.
证明 $\sqrt[3]{d}\,\omega$ 与 $\sqrt[3]{d}\,\omega^2$ 也是 $f(x)$ 的根,其中 $\omega$ 为三次单位根.
\end{example}
\exampleblank

\begin{solution}
由于 $\sqrt[3]{d}\notin\mathbb{Q}$,三次多项式
\[
  x^3-d
\]
在 $\mathbb{Q}[x]$ 中没有有理根,故不可约.它就是 $\sqrt[3]{d}$ 的最小多项式.
由 $f(\sqrt[3]{d})=0$ 可知 $x^3-d\mid f(x)$.因此 $x^3-d$ 的另外两个根
\[
  \sqrt[3]{d}\,\omega,
  \qquad
  \sqrt[3]{d}\,\omega^2
\]
也都是 $f$ 的根.
\end{solution}

%例题1.2.18
\begin{example}
设 $f(x)$ 是 $\mathbb{Q}$ 上的 $n$ 次不可约多项式.若 $f$ 的某个根的倒数也是 $f$ 的根,
证明:$f$ 的每一个根的倒数仍是 $f$ 的根.
\end{example}
\exampleblank

\begin{solution}
设 $f(x)=a_0x^n+a_1x^{n-1}+\cdots+a_n$,并令
\[
  f^*(x)=x^n f\left(\dfrac1x\right)
  =a_0+a_1x+\cdots+a_nx^n.
\]
若 $\alpha$ 与 $\alpha^{-1}$ 都是 $f$ 的根,则
\[
  f(\alpha)=0,
  \qquad
  f^*(\alpha)=\alpha^n f(\alpha^{-1})=0.
\]
故不可约多项式 $f$ 与 $f^*$ 有公共根,从而二者在 $\mathbb{Q}[x]$ 中相伴:
\[
  f^*=cf
\]
其中 $c\ne0$.因此若 $\beta$ 为 $f$ 的任意根,则
\[
  0=f^*(\beta^{-1})=\beta^{-n}f(\beta),
\]
等价地,$\beta^{-1}$ 也是 $f$ 的根.
\end{solution}


\newpage

\section{因式分解}

%例题1.3.1
\begin{example}[\markstar]
证明:实系数奇数次方程必有实数根.
\end{example}
\exampleblank

\begin{solution}
设
\[
  f(x)=a_{2m+1}x^{2m+1}+\cdots+a_0,
  \qquad a_{2m+1}\ne0.
\]
则当 $x\to+\infty$ 与 $x\to-\infty$ 时,首项决定 $f(x)$ 分别趋向异号的无穷大.
因此存在充分大的 $M>0$ 使 $f(M)f(-M)<0$.由连续函数介值定理,存在
$\xi\in(-M,M)$ 使 $f(\xi)=0$.
\end{solution}

%例题1.3.2
\begin{example}
证明:
\[
  (f^n,g^n)=(f,g)^n.
\]
\end{example}
\exampleblank

\begin{solution}
这与例题 1.2.12 相同.由不可约因式标准分解,最大公因式中每个不可约因式的指数等于两边指数的较小者,
故幂次整体乘以 $n$ 后仍有
\[
  (f^n,g^n)=(f,g)^n.
\]
\end{solution}

%例题1.3.3
\begin{example}
设 $f(x),g(x)$ 是数域 $F$ 上的多项式.证明:
\[
  f(x)\mid g(x)
\]
当且仅当对任意自然数 $n>1$ 都有
\[
  f^n(x)\mid g^n(x).
\]
\end{example}
\exampleblank

\begin{solution}
若 $f\mid g$,设 $g=fh$,则
\[
  g^n=f^nh^n,
\]
故 $f^n\mid g^n$.

反之,题设对任意 $n>1$ 成立,特别地对 $n=2$ 成立.由唯一分解定理,若某不可约因式 $p$ 在 $f$ 中指数为
$\alpha$,在 $g$ 中指数为 $\beta$,则 $f^2\mid g^2$ 给出 $2\alpha\le2\beta$,故 $\alpha\le\beta$.
于是 $f\mid g$.
\end{solution}

\subsection{单位分解}

\begin{theorem}
设复数
\[
  z_1=r_1(\cos\alpha+i\sin\alpha),
  \qquad
  z_2=r_2(\cos\beta+i\sin\beta),
\]
则
\[
  z_1z_2=r_1r_2\left[\cos(\alpha+\beta)+i\sin(\alpha+\beta)\right].
\]
\end{theorem}

\begin{corollary}[\markstar\ De Moivre 公式]
\[
  \left[r(\cos\theta+i\sin\theta)\right]^n
  =r^n\left(\cos n\theta+i\sin n\theta\right).
\]
\end{corollary}

%例题1.3.4
\begin{example}[\marktwostar]
求多项式 $x^n-1$ 分别在 $\mathbb{R}$,$\mathbb{C}$ 上的因式分解.
\end{example}
\exampleblank

\begin{solution}
在 $\mathbb{C}$ 上,$n$ 个根为
\[
  \omega_k=e^{\frac{2\pi ik}{n}},
  \qquad k=0,1,\ldots,n-1,
\]
故
\[
  x^n-1=\prod\limits_{k=0}^{n-1}
  \left(x-e^{\frac{2\pi ik}{n}}\right).
\]

在 $\mathbb{R}$ 上,将共轭根配对.若 $n$ 为奇数,
\[
  x^n-1=(x-1)
  \prod\limits_{k=1}^{\frac{n-1}{2}}
  \left(x^2-2x\cos\dfrac{2k\pi}{n}+1\right).
\]
若 $n$ 为偶数,
\[
  x^n-1=(x-1)(x+1)
  \prod\limits_{k=1}^{\frac n2-1}
  \left(x^2-2x\cos\dfrac{2k\pi}{n}+1\right).
\]
\end{solution}

%例题1.3.5
\begin{example}
求多项式 $x^n+1$ 分别在 $\mathbb{R}$,$\mathbb{C}$ 上的因式分解.
\end{example}
\exampleblank

\begin{solution}
在 $\mathbb{C}$ 上,根为
\[
  \zeta_k=e^{\frac{(2k+1)\pi i}{n}},
  \qquad k=0,1,\ldots,n-1,
\]
故
\[
  x^n+1=\prod\limits_{k=0}^{n-1}
  \left(x-e^{\frac{(2k+1)\pi i}{n}}\right).
\]
在 $\mathbb{R}$ 上将共轭根配对:若 $n$ 为偶数,
\[
  x^n+1=
  \prod\limits_{k=0}^{\frac n2-1}
  \left(x^2-2x\cos\dfrac{(2k+1)\pi}{n}+1\right).
\]
若 $n$ 为奇数,$-1$ 是实根,故
\[
  x^n+1=(x+1)
  \prod\limits_{k=0}^{\frac{n-3}{2}}
  \left(x^2-2x\cos\dfrac{(2k+1)\pi}{n}+1\right).
\]
\end{solution}

%例题1.3.6
\begin{example}
求多项式 $x^n-a^n$($a>0$)在 $\mathbb{R}$ 上的因式分解.
\end{example}
\exampleblank

\begin{solution}
令 $x=ay$,则
\[
  x^n-a^n=a^n(y^n-1).
\]
因此直接由上一题 $y^n-1$ 的实因式分解得到:若 $n$ 为奇数,
\[
  x^n-a^n=(x-a)
  \prod\limits_{k=1}^{\frac{n-1}{2}}
  \left(x^2-2ax\cos\dfrac{2k\pi}{n}+a^2\right),
\]
若 $n$ 为偶数,
\[
  x^n-a^n=(x-a)(x+a)
  \prod\limits_{k=1}^{\frac n2-1}
  \left(x^2-2ax\cos\dfrac{2k\pi}{n}+a^2\right).
\]
\end{solution}

%例题1.3.7
\begin{example}
证明:
\[
  \cos\dfrac{\pi}{2n+1}
  \cos\dfrac{2\pi}{2n+1}\cdots
  \cos\dfrac{n\pi}{2n+1}
  =\dfrac{1}{2^n}.
\]
\end{example}
\exampleblank

\begin{solution}
由
\[
  \sin 2x=2\sin x\cos x
\]
逐次使用可得
\[
  \displaystyle\prod\limits_{k=1}^{n}\cos\dfrac{k\pi}{2n+1}
  =2^{-n}
  \dfrac{\displaystyle\prod_{k=1}^{n}\sin\frac{2k\pi}{2n+1}}
        {\displaystyle\prod_{k=1}^{n}\sin\frac{k\pi}{2n+1}}.
\]
而模 $2n+1$ 看,$2,4,\ldots,2n$ 与 $1,2,\ldots,n$ 对应的正弦值只发生排列,
因为
\[
  \sin\dfrac{(2n+1-k)\pi}{2n+1}
  =\sin\dfrac{k\pi}{2n+1}.
\]
故分子分母相等,结论成立.
\end{solution}

\subsection{余数定理}

%例题1.3.8
\begin{example}
设 $m,n,p$ 都是非负整数,证明:
\[
  x^2+x+1\mid x^{3m}+x^{3n+1}+x^{3p+2}.
\]
\end{example}
\exampleblank

\begin{solution}
模 $x^2+x+1$ 有
\[
  x^3-1=(x-1)(x^2+x+1)\equiv0,
\]
故 $x^3\equiv1$.于是
\[
  x^{3m}+x^{3n+1}+x^{3p+2}
  \equiv1+x+x^2\equiv0.
\]
\end{solution}

%例题1.3.9
\begin{example}
已知
\[
  x^2+x+1\mid f_1(x^3)+xf_2(x^3),
\]
证明:$x-1\mid f_1(x)$ 且 $x-1\mid f_2(x)$.
\end{example}
\exampleblank

\begin{solution}
令 $\omega,\omega^2$ 为 $x^2+x+1$ 的两个根,则 $\omega^3=1$ 且 $\omega\ne\omega^2$.
代入题设得到
\[
\begin{cases}
  f_1(1)+\omega f_2(1)=0,\\
  f_1(1)+\omega^2 f_2(1)=0.
\end{cases}
\]
相减得 $(\omega-\omega^2)f_2(1)=0$,故 $f_2(1)=0$,进而 $f_1(1)=0$.
由余数定理,$x-1$ 同时整除 $f_1,f_2$.
\end{solution}

%例题1.3.10
\begin{example}
假设
\[
  f_0(x^5)+xf_1(x^{10})+x^2f_2(x^{15})+x^3f_3(x^{20})
\]
能被 $x^4+x^3+x^2+x+1$ 整除.证明:
\[
  x-1\mid f_i(x),
  \qquad i=0,1,2,3.
\]
\end{example}
\exampleblank

\begin{solution}
设 $\omega$ 为任意非 $1$ 的五次单位根,则 $\omega^5=1$.代入题设多项式得
\[
  f_0(1)+\omega f_1(1)+\omega^2f_2(1)+\omega^3f_3(1)=0.
\]
四个非 $1$ 的五次单位根都是三次多项式
\[
  F(t)=f_0(1)+t f_1(1)+t^2f_2(1)+t^3f_3(1)
\]
的根.次数不超过 $3$ 的非零多项式不可能有四个不同根,故 $F\equiv0$.
于是 $f_i(1)=0$($i=0,1,2,3$),即 $x-1\mid f_i(x)$.
\end{solution}

%例题1.3.11
\begin{example}
设 $f(x),g(x),h(x)$ 满足
\[
\begin{cases}
  (x^2+1)h(x)+(x+1)f(x)+(x-2)g(x)=0,\\
  (x^2+1)h(x)+(x-1)f(x)+(x+2)g(x)=0.
\end{cases}
\]
证明:$x^2+1$ 是 $f(x)$ 和 $g(x)$ 的公因式.
\end{example}
\exampleblank

\begin{solution}
两式相减得
\[
  2f-4g=0,
\]
故 $f=2g$.代回第一式:
\[
  (x^2+1)h+3xg=0.
\]
由于
\[
  (x^2+1,x)=1,
\]
故 $x^2+1\mid g$,于是也有 $x^2+1\mid f$.
\end{solution}

%例题1.3.12
\begin{example}
证明:
\[
  x(x+1)(2x+1)
  \mid
  (x+1)^{2n}-x^{2n}-2x-1.
\]
\end{example}
\exampleblank

\begin{solution}
记右端为 $P(x)$.直接计算
\[
  P(0)=0,
  \qquad
  P(-1)=0,
  \qquad
  P\left(-\dfrac12\right)=0.
\]
故 $x$,$x+1$,$2x+1$ 均整除 $P(x)$.三者两两互素,因此
\[
  x(x+1)(2x+1)\mid P(x).
\]
\end{solution}

%例题1.3.13
\begin{example}
设 $f(x)$ 是一个 $n$ 次多项式,且当 $k=0,1,2,\ldots,n$ 时,
\[
  f(k)=k^{k+1}.
\]
求 $f(n+1)$.
\end{example}
\exampleblank

\begin{solution}
由于 $\deg f\le n$,其 $(n+1)$ 阶前向差分恒为零:
\[
  \displaystyle\sum\limits_{k=0}^{n+1}
  (-1)^{n+1-k}C_{n+1}^{k}f(k)=0.
\]
因此
\[
  \boxed{
  f(n+1)=
  \sum\limits_{k=0}^{n}
  (-1)^{n-k}C_{n+1}^{k}k^{k+1}}
\]
(其中 $k=0$ 项按 $0^{1}=0$ 处理).
\end{solution}

\subsection{重因式}

%例题1.3.14
\begin{example}
判断下列多项式有无重因式:
\begin{enumerate}
  \item $f(x)=x^5-10x^3-20x^2-15x-4$;
  \item $f(x)=4x^3-4x^2-7x-2$;
  \item $f(x)=x^3+x+1$.
\end{enumerate}
\end{example}
\exampleblank

\begin{solution}
直接分解:
\[
\begin{aligned}
  x^5-10x^3-20x^2-15x-4&=(x-4)(x+1)^4,\\
  4x^3-4x^2-7x-2&=(x-2)(2x+1)^2.
\end{aligned}
\]
故前两式分别有重因式 $x+1$ 与 $2x+1$.第三式
\[
  (x^3+x+1,3x^2+1)=1,
\]
故无重因式.
\end{solution}

%例题1.3.15
\begin{example}
已知
\[
  f(x)=x^3+6x^2+3px+8,
\]
确定 $p$ 的值,使 $f(x)$ 有重根,并求其根.
\end{example}
\exampleblank

\begin{solution}
三次多项式有重根当且仅当判别式为零.计算得
\[
  \Delta(f)=-108(p-4)^2(p+5).
\]
故
\[
  p=4\quad\text{或}\quad p=-5.
\]
当 $p=4$ 时,
\[
  f(x)=(x+2)^3;
\]
当 $p=-5$ 时,
\[
  f(x)=(x-1)^2(x+8).
\]
\end{solution}

%例题1.3.16
\begin{example}[\markstar]
证明:
\[
  f(x)=1+\dfrac{x}{1!}+\dfrac{x^2}{2!}+\cdots+\dfrac{x^n}{n!}
\]
无重根.
\end{example}
\exampleblank

\begin{solution}
记
\[
  f_n(x)=1+\dfrac{x}{1!}+\cdots+\dfrac{x^n}{n!}.
\]
则
\[
  f_n'(x)=f_{n-1}(x),
  \qquad
  f_n(x)-f_{n-1}(x)=\dfrac{x^n}{n!}.
\]
若 $\alpha$ 是重根,则
\[
  f_n(\alpha)=f_n'(\alpha)=f_{n-1}(\alpha)=0,
\]
从而 $\alpha^n=0$,故 $\alpha=0$.但 $f_n(0)=1$,矛盾.因此无重根.
\end{solution}

\begin{proposition}[\markstar]
\[
  \dfrac{f(x)}{(f(x),f'(x))}
\]
是一个没有重因式的多项式,并且它的不可约因式与 $f(x)$ 的不可约因式相同.
\end{proposition}

%例题1.3.17
\begin{example}
设 $\deg f(x)=n\ge1$.若
\[
  f'(x)\mid f(x),
\]
证明:$f(x)$ 有一个 $n$ 重根.
\end{example}
\exampleblank

\begin{solution}
因为 $\deg f'=n-1$,存在一次多项式 $ax+b$ 使
\[
  f=(ax+b)f'.
\]
比较最高次项可知 $a=\dfrac1n$.令 $\alpha=-b/a$,则
\[
  f(x)=\dfrac{x-\alpha}{n}f'(x).
\]
在 $x\ne\alpha$ 处可写成
\[
  \dfrac{f'(x)}{f(x)}=\dfrac{n}{x-\alpha}.
\]
积分或比较多项式即可得到
\[
  f(x)=C(x-\alpha)^n,
\]
故 $\alpha$ 是 $n$ 重根.
\end{solution}


\newpage

\section{有理,整系数多项式}

%例题1.4.1
\begin{example}
求下列多项式的有理根:
\begin{enumerate}
  \item $x^3-6x^2+15x-14$;
  \item $4x^4-7x^2-5x-1$;
  \item $x^5-12x^3+36x+12$.
\end{enumerate}
\end{example}
\exampleblank

\begin{solution}
由有理根定理逐一检验候选值:
\begin{enumerate}
  \item
  \[
    x^3-6x^2+15x-14=(x-2)(x^2-4x+7),
  \]
  故唯一有理根为 $2$;
  \item
  \[
    4x^4-7x^2-5x-1=(2x+1)^2(x^2-x-1),
  \]
  故唯一有理根为 $-\dfrac12$,重数为 $2$;
  \item 候选有理根为 $\pm1,\pm2,\pm3,\pm4,\pm6,\pm12$,逐一代入均不为零,故无有理根.
\end{enumerate}
\end{solution}

%例题1.4.2
\begin{example}
判断下列多项式在 $\mathbb{Q}$ 上是否可约:
\begin{enumerate}
  \item $x^4-8x^3+12x^2+2$;
  \item $x^p+px+1$,其中 $p$ 为奇素数;
  \item $x^{p-1}+x^{p-2}+\cdots+x+1$,其中 $p$ 为素数.
\end{enumerate}
\end{example}
\exampleblank

\begin{solution}
三者均在 $\mathbb{Q}$ 上不可约.

(1) 模 $3$ 化简得到
\[
  x^4+x^3-1\in\mathbb{F}_3[x].
\]
它在 $\mathbb{F}_3$ 中无一次因子;进一步检验唯一可能的首一不可约二次因子可知亦无二次因子,
故在 $\mathbb{F}_3[x]$ 中不可约.由模素数判别法,原多项式在 $\mathbb{Q}[x]$ 中不可约.

(2) 作平移 $x=y-1$:
\[
  (y-1)^p+p(y-1)+1.
\]
由于 $p$ 为奇素数,除首项外各系数均被 $p$ 整除,而常数项为 $-p$,不被 $p^2$ 整除.
由 Eisenstein 判别法知其不可约,故原多项式不可约.

(3) 这是素数阶分圆多项式
\[
  \Phi_p(x)=\dfrac{x^p-1}{x-1}.
\]
作平移 $x=y+1$ 后,除首项外各项系数均被 $p$ 整除,而常数项为 $p$ 且不被 $p^2$ 整除.
由 Eisenstein 判别法得不可约.
\end{solution}

%例题1.4.3
\begin{example}
设 $f(x)$ 是次数大于 $0$ 的首一整系数多项式.若 $f(0),f(1)$ 都是奇数,证明:$f(x)$ 没有有理根.
\end{example}
\exampleblank

\begin{solution}
首一整系数多项式的有理根必为整数.若整数根 $m$ 为偶数,则
\[
  f(m)\equiv f(0)\equiv1\pmod2,
\]
不可能为 $0$;若 $m$ 为奇数,则
\[
  f(m)\equiv f(1)\equiv1\pmod2,
\]
同样不可能为 $0$.故无有理根.
\end{solution}

%例题1.4.4
\begin{example}[\markstar]
设 $f(x)$ 是整系数多项式,且
\[
  f(0),f(1),\ldots,f(n-1)
\]
都不能被 $n$ 整除.证明:$f(x)$ 没有整数根.
\end{example}
\exampleblank

\begin{solution}
若整数 $m$ 是 $f$ 的根,取 $r\in\{0,1,\ldots,n-1\}$ 使 $m\equiv r\pmod n$.
整系数多项式保持同余,故
\[
  f(m)\equiv f(r)\pmod n.
\]
由于 $f(m)=0$,得到 $n\mid f(r)$,与题设矛盾.
\end{solution}

%例题1.4.5
\begin{example}
设 $a_1,a_2,\ldots,a_n$ 为互不相同的整数,令
\[
  P(x)=\prod\limits_{i=1}^{n}(x-a_i),
  \qquad
  f(x)=P(x)-1.
\]
\begin{enumerate}
  \item 证明 $f(x)$ 在 $\mathbb{Q}$ 上不可约;
  \item 对整数 $t\ne-1$,讨论
  \[
    g(x)=P(x)+t
  \]
  在 $\mathbb{Q}$ 上是否一定不可约.
\end{enumerate}
\end{example}
\exampleblank

\begin{solution}

(1) 反设 $f=gh$ 在 $\mathbb{Z}[x]$ 中有非平凡分解,并取
\[
  0<\deg g\le\deg h.
\]
由 $f(a_i)=-1$,知
\[
  g(a_i)=\pm1.
\]
于是 $g(x)^2-1$ 在 $a_1,\ldots,a_n$ 处均为零,故
\[
  P(x)\mid g(x)^2-1.
\]
又 $2\deg g\le n$.若 $2\deg g<n$,则次数矛盾;若 $2\deg g=n$,由首项系数比较只能有
\[
  g(x)^2-1=P(x).
\]
于是
\[
  f(x)=P(x)-1=g(x)^2-2.
\]
但 $g\mid f$ 将推出 $g\mid2$,与 $\deg g>0$ 矛盾.故 $f$ 不可约.

(2) 对 $t\ne-1$ 不能保证不可约.例如取 $t=0$,则
\[
  g(x)=P(x)=\prod\limits_{i=1}^{n}(x-a_i)
\]
显然可约(当 $n>1$).因此题设中的 $t\ne-1$ 并不足以推出统一的不可约性结论.
\end{solution}

%例题1.4.6
\begin{example}
设 $a_1,a_2,\ldots,a_n$ 是互不相同的整数.证明
\[
  f(x)=\prod\limits_{i=1}^{n}(x-a_i)+1
\]
在 $\mathbb{Q}$ 上不可约,或者是某一有理系数多项式的平方.
\end{example}
\exampleblank

\begin{solution}
记
\[
  P(x)=\prod\limits_{i=1}^{n}(x-a_i).
\]
若 $f=P+1$ 可约,由 Gauss 引理可在 $\mathbb{Z}[x]$ 中写成 $f=gh$,并取
$0<\deg g\le\deg h$.由
\[
  f(a_i)=1
\]
知 $g(a_i)=h(a_i)=\pm1$,因而 $P\mid g^2-1$.与上一题同样的次数比较表明
$2\deg g<n$ 不可能,故 $2\deg g=n$,并由首项系数比较得到
\[
  P=g^2-1.
\]
于是
\[
  f=P+1=g^2.
\]
因此若 $f$ 可约,它必为一个有理(事实上整系数)多项式的平方.
\end{solution}

%例题1.4.7
\begin{example}
设 $a_1,a_2,\ldots,a_n$ 是 $n$ 个互不相同的整数,证明
\[
  f(x)=\prod\limits_{i=1}^{n}(x-a_i)^2+1
\]
在 $\mathbb{Q}$ 上不可约.
\end{example}
\exampleblank

\begin{solution}
记
\[
  P(x)=\prod\limits_{i=1}^{n}(x-a_i),
  \qquad
  f=P^2+1.
\]
若 $f=gh$ 在 $\mathbb{Z}[x]$ 中非平凡分解,则 $f(x)>0$ 对所有实数 $x$ 成立,故 $g,h$ 均无实根,
从而各自在 $\mathbb{R}$ 上不变号.又
\[
  g(a_i)h(a_i)=1,
\]
所以对所有 $i$,$g(a_i)$ 必恒为同一个 $1$ 或恒为同一个 $-1$.

不妨 $0<\deg g\le n$.若 $\deg g<n$,则 $g(x)\mp1$ 有 $n$ 个不同整数根,迫使 $g\equiv\pm1$,矛盾.
故 $\deg g=n$,同理 $\deg h=n$.同时 $g(a_i)=h(a_i)$,故 $P\mid g-h$.
取两个因子的首项系数同号并首一化后,$\deg(g-h)<n$,于是 $g=h$.这将给出
\[
  P^2+1=g^2,
\]
即
\[
  (g-P)(g+P)=1,
\]
不可能.故 $f$ 在 $\mathbb{Q}$ 上不可约.
\end{solution}

%例题1.4.8
\begin{example}
设
\[
  f(x)=\prod\limits_{i=1}^{2025}(x-i)^2+2026.
\]
判断 $f(x)$ 在 $\mathbb{Q}$ 上是否可约,并证明结论.
\end{example}
\exampleblank

\begin{solution}
记
\[
  P(x)=\prod\limits_{i=1}^{2025}(x-i),
  \qquad
  F(x)=P^2(x)+2026.
\]
反设 $F$ 在 $\mathbb{Q}[x]$ 中可约.由于 $F$ 为首一整系数多项式,由 Gauss 引理可写成
\[
  F(x)=g(x)h(x),
\]
其中 $g,h\in\mathbb{Z}[x]$ 均首一,且
\[
  0<\deg g\leq\deg h.
\]
于是
\[
  \deg g\leq2025.
\]
又因为
\[
  F(x)=P^2(x)+2026>0
\]
对所有实数 $x$ 成立,所以 $g,h$ 都没有实根.二者均首一,因此在实轴上恒正.
特别地,对 $i=1,2,\ldots,2025$,
\[
  g(i)h(i)=F(i)=2026=2\times1013,
\]
从而
\[
  g(i)\in\{1,2,1013,2026\}.
\]

对任意整数 $u,v$,有
\[
  u-v\mid g(u)-g(v).
\]
注意集合 $\{1,2,1013,2026\}$ 中两个不同元素之差的绝对值只能是
\[
  1,
  \ 1011,
  \ 1012,
  \ 1013,
  \ 2024,
  \ 2025,
\]
这些数没有一个能被 $6$ 或 $7$ 整除.因此,只要指标仍落在
$\{1,2,\ldots,2025\}$ 内,就有
\[
  g(i+6)=g(i),
  \qquad
  g(i+7)=g(i).
\]
于是对 $1\leq i\leq2018$,
\[
  g(i)=g(i+7)=g(i+1),
\]
再利用步长 $6$ 可把最后几个整数也与前面连接起来,故
\[
  g(1)=g(2)=\cdots=g(2025)=c
\]
对某个 $c\in\{1,2,1013,2026\}$ 成立.

因此 $g(x)-c$ 有 $2025$ 个不同整数根.由于 $g$ 非常数且
$\deg g\leq2025$,只能有
\[
  \deg g=2025,
  \qquad
  g(x)=P(x)+c.
\]
于是 $h$ 也为 $2025$ 次首一多项式.由同样的论证,存在整数 $d>0$ 使
\[
  h(x)=P(x)+d.
\]
比较
\[
  (P+c)(P+d)=P^2+2026
\]
的 $P$ 项与常数项,得到
\[
  c+d=0,
  \qquad
  cd=2026,
\]
矛盾.

故
\[
  \boxed{F(x)\text{ 在 }\mathbb{Q}[x]\text{ 上不可约}.}
\]
\end{solution}

%例题1.4.9
\begin{example}
设 $f(x)$ 是整系数多项式,且
\[
  f(1)=f(2)=f(3)=p,
\]
其中 $p$ 为素数.证明:不存在整数 $m$ 使 $f(m)=2p$.
\end{example}
\exampleblank

\begin{solution}
若存在整数 $m$ 使 $f(m)=2p$,则对 $i=1,2,3$,整系数多项式的差值性质给出
\[
  m-i\mid f(m)-f(i)=p.
\]
故每个 $m-i$ 都只能取 $\pm1,\pm p$ 的某个值.
但 $m-1,m-2,m-3$ 是三个连续整数,不可能全部属于集合
$\{\pm1,\pm p\}$.矛盾.
\end{solution}

%例题1.4.10
\begin{example}[\markstar]
设 $f(x)$ 是实系数多项式,且对任意实数 $c$ 都有 $f(c)\ge0$.证明:存在实系数多项式 $g,h$ 使
\[
  f(x)=g^2(x)+h^2(x).
\]
\end{example}
\exampleblank

\begin{solution}
将 $f$ 在 $\mathbb{R}[x]$ 中分解.由于 $f(x)\ge0$,每个实根的重数都是偶数,故实一次因子的乘积可并入一个平方.
其余不可约二次因子均形如
\[
  (x-a)^2+b^2,
  \qquad b\ne0,
\]
本身就是两个实系数多项式平方之和.
再反复使用恒等式
\[
  (u^2+v^2)(s^2+t^2)
  =(us-vt)^2+(ut+vs)^2,
\]
即可将全部因子的乘积写成两个实系数多项式平方之和.
\end{solution}

%例题1.4.11
\begin{example}
若
\[
  f(x)\mid f(x^m),
  \qquad m>1,
\]
证明:$f(x)$ 的根只能是 $0$ 或单位根.
\end{example}
\exampleblank

\begin{solution}
设 $\alpha$ 为 $f$ 的根.由 $f\mid f(x^m)$ 得
\[
  f(\alpha^m)=0.
\]
继续迭代可知
\[
  \alpha,\alpha^m,\alpha^{m^2},\ldots
\]
全是 $f$ 的根.若 $\alpha=0$,结论成立;若 $\alpha\ne0$,由于 $f$ 只有有限多个根,
必有 $r<s$ 使
\[
  \alpha^{m^r}=\alpha^{m^s}.
\]
于是
\[
  \alpha^{m^r(m^{s-r}-1)}=1.
\]
故 $\alpha$ 是单位根.
\end{solution}


\newpage

\section{多元多项式}

\begin{theorem}[\marktwostar\ Vieta 定理]
对于数域 $K$ 上任意 $n$ 次多项式
\[
  f(x)=a_0x^n+a_1x^{n-1}+\cdots+a_{n-1}x+a_n,
  \qquad a_0\ne0,
\]
设 $x_1,\ldots,x_n$ 为其在复数域中的全部根(按重数计),则
\[
  \sigma_i
  =\sum\limits_{1\le k_1<\cdots<k_i\le n}
  x_{k_1}\cdots x_{k_i}
  =(-1)^i\dfrac{a_i}{a_0},
  \qquad i=1,\ldots,n.
\]
\end{theorem}

\begin{theorem}[\markstar\ Newton 公式]
令
\[
  s_k=x_1^k+x_2^k+\cdots+x_n^k.
\]
当 $1\le k\le n$ 时,
\[
  s_k-\sigma_1s_{k-1}+\sigma_2s_{k-2}-\cdots
  +(-1)^{k-1}\sigma_{k-1}s_1+(-1)^k k\sigma_k=0.
\]
当 $k>n$ 时,
\[
  s_k-\sigma_1s_{k-1}+\sigma_2s_{k-2}-\cdots
  +(-1)^n\sigma_ns_{k-n}=0.
\]
\end{theorem}

\begin{theorem}[\markstar]
设 $f(x_1,\ldots,x_n)$ 是数域 $K$ 上的对称多项式,则存在唯一的
$g(y_1,\ldots,y_n)\in K[y_1,\ldots,y_n]$ 使
\[
  f(x_1,\ldots,x_n)=g(\sigma_1,\ldots,\sigma_n).
\]
\end{theorem}

%例题1.5.1
\begin{example}
按字典排列法整理下列多项式:
\begin{enumerate}
  \item $x_1^2x_2+x_1^3x_3+3x_1x_2x_3-4x_2x_3^5$;
  \item $(x_1^3-2x_3+x_1x_2x_3)(x_1^2-x_3^2-x_1x_2)$.
\end{enumerate}
\end{example}
\exampleblank

\begin{solution}
按 $x_1>x_2>x_3$ 的字典序:
\begin{enumerate}
  \item
  \[
    x_1^3x_3+x_1^2x_2+3x_1x_2x_3-4x_2x_3^5.
  \]
  \item 展开并排序得
  \[
  \begin{aligned}
    {}&x_1^5-x_1^4x_2+x_1^3x_2x_3-x_1^3x_3^2
    -x_1^2x_2^2x_3-2x_1^2x_3\\
    &\qquad -x_1x_2x_3^3+2x_1x_2x_3+2x_3^3.
  \end{aligned}
  \]
\end{enumerate}
\end{solution}

%例题1.5.2
\begin{example}
设 $f(x_1,\ldots,x_n),g(x_1,\ldots,x_n)$ 是数域 $K$ 上的多项式,且 $g$ 不是零多项式.
若对一切满足 $g(a_1,\ldots,a_n)\ne0$ 的 $(a_1,\ldots,a_n)\in K^n$ 都有
\[
  f(a_1,\ldots,a_n)=0,
\]
证明 $f$ 是零多项式.
\end{example}
\exampleblank

\begin{solution}
数域 $K$ 为无限域.对变量个数作归纳.$n=1$ 时,非零一元多项式只有有限个根,
而 $g$ 的零点也有限,因此不可能在所有 $g(a)\ne0$ 的点上均有 $f(a)=0$,除非 $f\equiv0$.

设结论对 $n-1$ 个变量成立.将 $f,g$ 视为关于 $x_n$ 的多项式,其系数属于
$K[x_1,\ldots,x_{n-1}]$.若 $f\not\equiv0$,可选取
$(a_1,\ldots,a_{n-1})$ 使得 $f$ 的某个非零系数与 $g$ 的某个非零系数均不为零.
于是得到两个非零的一元多项式
\[
  f(a_1,\ldots,a_{n-1},x_n),
  \qquad
  g(a_1,\ldots,a_{n-1},x_n).
\]
在无限域中可选 $a_n$ 使后者不为零且前者也不为零,与题设矛盾.因此 $f\equiv0$.
\end{solution}

%例题1.5.3
\begin{example}
用初等对称多项式 $\sigma_1,\sigma_2,\sigma_3$ 表示:
\begin{enumerate}
  \item $(x_1+x_2+x_1x_2)(x_2+x_3+x_2x_3)(x_1+x_3+x_1x_3)$;
  \item $(x_1^2+x_2x_3)(x_2^2+x_1x_3)(x_3^2+x_1x_2)$;
  \item $x_1^3+x_2^3+x_3^3-3x_1x_2x_3$.
\end{enumerate}
\end{example}
\exampleblank

\begin{solution}
整理得:
\begin{enumerate}
  \item
  \[
    \sigma_1\sigma_2+\sigma_1\sigma_3+\sigma_2^2
    +2\sigma_2\sigma_3+\sigma_3^2-\sigma_3;
  \]
  \item
  \[
    \sigma_1^3\sigma_3-6\sigma_1\sigma_2\sigma_3
    +\sigma_2^3+8\sigma_3^2;
  \]
  \item
  \[
    \sigma_1^3-3\sigma_1\sigma_2.
  \]
\end{enumerate}
\end{solution}

%例题1.5.4
\begin{example}
设方程
\[
  x^3+px^2+qx+r=0
\]
的三根成等差数列,证明
\[
  2p^3-9pq+27r=0.
\]
\end{example}
\exampleblank

\begin{solution}
设三根为 $m-d,m,m+d$.由 Vieta 定理
\[
  3m=-p,
\]
且
\[
  q=(m-d)m+(m-d)(m+d)+m(m+d)=3m^2-d^2,
\]
\[
  -r=(m-d)m(m+d)=m(m^2-d^2).
\]
由 $d^2=3m^2-q$ 得
\[
  r=-m(q-2m^2).
\]
再代入 $m=-p/3$,整理即得
\[
  2p^3-9pq+27r=0.
\]
\end{solution}

%例题1.5.5
\begin{example}
若方程
\[
  x^3+px^2+qx+r=0
\]
的三根都是实数,证明
\[
  p^2\ge3q.
\]
\end{example}
\exampleblank

\begin{solution}
设三实根为 $x_1,x_2,x_3$.由 Vieta 定理
\[
  x_1+x_2+x_3=-p,
  \qquad
  x_1x_2+x_1x_3+x_2x_3=q.
\]
又
\[
  (x_1+x_2+x_3)^2-3(x_1x_2+x_1x_3+x_2x_3)
  =\dfrac12\sum_{i<j}(x_i-x_j)^2\ge0.
\]
故 $p^2\ge3q$.
\end{solution}

\begin{proposition}[\markstar\ Cardan 公式]
方程
\[
  x^3+ax^2+bx+c=0
\]
通过变换 $x=y-\dfrac a3$ 可化为
\[
  y^3+py+q=0.
\]
令
\[
  \Delta=\dfrac{q^2}{4}+\dfrac{p^3}{27},
  \qquad
  u=\sqrt[3]{-\dfrac q2+\sqrt\Delta},
  \qquad
  v=\sqrt[3]{-\dfrac q2-\sqrt\Delta},
\]
并令 $\omega$ 为三次单位根,则三个根可写为
\[
  u+v,
  \qquad
  \omega u+\omega^2v,
  \qquad
  \omega^2u+\omega v,
\]
再反代 $x=y-\dfrac a3$ 即得原方程的根.
\end{proposition}

%例题1.5.6
\begin{example}[\markstar]
证明:方程
\[
  x^3+px+q=0
\]
有重根的充要条件是
\[
  4p^3+27q^2=0.
\]
\end{example}
\exampleblank

\begin{solution}
有重根当且仅当 $f(x)=x^3+px+q$ 与
\[
  f'(x)=3x^2+p
\]
有公共根 $\alpha$.于是
\[
  3\alpha^2+p=0,
  \qquad
  \alpha^3+p\alpha+q=0.
\]
第一式给出 $p=-3\alpha^2$,代入第二式得 $q=2\alpha^3$,从而
\[
  4p^3+27q^2
  =4(-3\alpha^2)^3+27(2\alpha^3)^2=0.
\]
反之,若 $4p^3+27q^2=0$,即该三次多项式判别式为零,故必有重根.
\end{solution}

%例题1.5.7
\begin{example}
已知
\[
  x^3+px^2+qx+r=0
\]
的三个根为 $\alpha_1,\alpha_2,\alpha_3$.求一个以
$\alpha_1^3,\alpha_2^3,\alpha_3^3$ 为根的三次方程.
\end{example}
\exampleblank

\begin{solution}
由 Vieta 定理
\[
  e_1=\alpha_1+\alpha_2+\alpha_3=-p,
  \quad
  e_2=q,
  \quad
  e_3=-r.
\]
令 $\beta_i=\alpha_i^3$.则
\[
\begin{aligned}
  \beta_1+\beta_2+\beta_3
  &=e_1^3-3e_1e_2+3e_3
   =-p^3+3pq-3r,\\
  \beta_1\beta_2+\beta_1\beta_3+\beta_2\beta_3
  &=q^3-3pqr+3r^2,\\
  \beta_1\beta_2\beta_3&=-r^3.
\end{aligned}
\]
故所求方程为
\[
\boxed{
  y^3+(p^3-3pq+3r)y^2
  +(q^3-3pqr+3r^2)y+r^3=0 }.
\]
\end{solution}

%例题1.5.8
\begin{example}
解方程组
\[
\begin{cases}
  x+y+z+w=10,\\
  x^2+y^2+z^2+w^2=30,\\
  x^3+y^3+z^3+w^3=100,\\
  xyzw=24.
\end{cases}
\]
\end{example}
\exampleblank

\begin{solution}
设 $e_k$ 为 $x,y,z,w$ 的初等对称多项式.由 Newton 公式,
\[
  e_1=10,
\]
\[
  e_2=\dfrac{e_1^2-s_2}{2}
  =\dfrac{100-30}{2}=35,
\]
\[
  e_3=\dfrac{s_3-e_1s_2+e_2s_1}{3}
  =\dfrac{100-300+350}{3}=50,
\]
且 $e_4=xyzw=24$.故 $x,y,z,w$ 是多项式
\[
  t^4-10t^3+35t^2-50t+24
  =(t-1)(t-2)(t-3)(t-4)
\]
的四个根.因此
\[
  \{x,y,z,w\}=\{1,2,3,4\},
\]
即全部解由 $(1,2,3,4)$ 的任意排列给出.
\end{solution}

%例题1.5.9
\begin{example}
设 $x_1,x_2,x_3$ 是方程
\[
  x^3-6x^2+ax+a=0
\]
的三个根.求使
\[
  (x_1-1)^3+(x_2-2)^3+(x_3-3)^3=0
\]
成立的所有实数 $a$,并求相应的 $x_1,x_2,x_3$.
\end{example}
\exampleblank

\begin{solution}
由 Vieta 关系
\[
  x_1+x_2+x_3=6,
  \qquad
  x_1x_2+x_1x_3+x_2x_3=a,
  \qquad
  x_1x_2x_3=-a.
\]
与附加条件联立消元,可得
\[
  (2a-5)(3a-16)(4a-27)=0.
\]
因此
\[
  a=\dfrac52,\qquad \dfrac{16}{3},\qquad \dfrac{27}{4}.
\]
相应的有序根为:
\[
\begin{array}{c|c}
 a & (x_1,x_2,x_3)\\ \hline
 \dfrac52 & \left(1,\dfrac{5\pm\sqrt{35}}2,\dfrac{5\mp\sqrt{35}}2\right)\\[1.2ex]
 \dfrac{16}{3} & \left(2\pm\dfrac{2\sqrt{15}}3,2,2\mp\dfrac{2\sqrt{15}}3\right)\\[1.2ex]
 \dfrac{27}{4} & \left(\dfrac{3\pm3\sqrt2}{2},\dfrac{3\mp3\sqrt2}{2},3\right)
\end{array}
\]
其中每一行的 $\pm$ 与 $\mp$ 同时取相反符号.
\end{solution}

%例题1.5.10
\begin{example}
计算多项式
\[
  x^9+x^7+x^3+x^2+x-1
\]
所有根的平方和,立方和及九次方之和.
\end{example}
\exampleblank

\begin{solution}
记根的 $k$ 次幂和为 $s_k$.由 Newton 公式递推可得
\[
  s_1=0,
  \qquad
  s_2=-2,
  \qquad
  s_3=0,
\]
并继续递推到九次得到
\[
  s_9=18.
\]
故所求为
\[
  \boxed{s_2=-2,\qquad s_3=0,\qquad s_9=18}.
\]
\end{solution}

%例题1.5.11
\begin{example}
解方程组
\[
\begin{cases}
  x_1+x_2+x_3+x_4=4,\\
  x_1^2+x_2^2+x_3^2+x_4^2=4,\\
  x_1^3+x_2^3+x_3^3+x_4^3=4,\\
  x_1^4+x_2^4+x_3^4+x_4^4=4.
\end{cases}
\]
\end{example}
\exampleblank

\begin{solution}
由 Newton 公式依次得到
\[
  e_1=4,
  \qquad e_2=6,
  \qquad e_3=4,
  \qquad e_4=1.
\]
因此 $x_1,x_2,x_3,x_4$ 都是
\[
  t^4-4t^3+6t^2-4t+1=(t-1)^4
\]
的根,故唯一解为
\[
  \boxed{x_1=x_2=x_3=x_4=1}.
\]
\end{solution}

%例题1.5.12
\begin{example}[\markstar]
若对任意正整数 $k$ 都有
\[
  \tr(A^k)=0,
\]
证明:$A$ 的特征值均为零.
\end{example}
\exampleblank

\begin{solution}
设 $A$ 为 $n$ 阶方阵,特征值为 $\lambda_1,\ldots,\lambda_n$(按代数重数计).则
\[
  \tr(A^k)=\sum\limits_{i=1}^{n}\lambda_i^k=0,
  \qquad k=1,2,\ldots,n.
\]
由 Newton 公式,前 $n$ 个幂和全为零推出全部初等对称多项式
\[
  \sigma_1=\sigma_2=\cdots=\sigma_n=0.
\]
故 $A$ 的特征多项式为
\[
  \lambda^n,
\]
从而所有特征值均为零.
\end{solution}

%例题1.5.13
\begin{example}[\markstar]
证明:对 $1\le k\le n$,有
\[
  s_k=
  \begin{vmatrix}
    \sigma_1&1&0&\cdots&0\\
    2\sigma_2&\sigma_1&1&\cdots&0\\
    3\sigma_3&\sigma_2&\sigma_1&\ddots&\vdots\\
    \vdots&\vdots&\ddots&\ddots&1\\
    k\sigma_k&\sigma_{k-1}&\cdots&\sigma_2&\sigma_1
  \end{vmatrix},
\]
以及
\[
  \sigma_k=\dfrac1{k!}
  \begin{vmatrix}
    s_1&1&0&\cdots&0\\
    s_2&s_1&2&\cdots&0\\
    s_3&s_2&s_1&\ddots&\vdots\\
    \vdots&\vdots&\ddots&\ddots&k-1\\
    s_k&s_{k-1}&\cdots&s_2&s_1
  \end{vmatrix}.
\]
\end{example}
\exampleblank

\begin{solution}
Newton 公式对 $j=1,\ldots,k$ 构成关于 $s_1,\ldots,s_k$ 的三角线性方程组:
\[
  s_j-\sigma_1s_{j-1}+\cdots+(-1)^j j\sigma_j=0.
\]
将其按 $s_1,\ldots,s_k$ 排列后,用 Cramer 法则求 $s_k$,即得到第一个行列式公式.

反过来,把同一组 Newton 公式视为关于 $\sigma_1,\ldots,\sigma_k$ 的三角线性方程组,
其系数矩阵行列式为 $k!$.再次应用 Cramer 法则便得到第二个公式.
\end{solution}


\newpage

\section{结式和判别式}

\begin{lemma}
设
\[
  (f(x),g(x))=d(x).
\]
则 $d(x)\ne1$ 的充要条件是存在 $K[x]$ 中非零多项式 $u(x),v(x)$,满足
\[
  f(x)u(x)=g(x)v(x),
\]
且
\[
  \deg u<\deg g,
  \qquad
  \deg v<\deg f.
\]
\end{lemma}

设
\[
\begin{aligned}
  f(x)&=a_0x^n+a_1x^{n-1}+\cdots+a_n,\\
  g(x)&=b_0x^m+b_1x^{m-1}+\cdots+b_m.
\end{aligned}
\]
由上引理,将
\[
  u(x)=x_0x^{m-1}+x_1x^{m-2}+\cdots+x_{m-1},
\]
\[
  v(x)=y_0x^{n-1}+y_1x^{n-2}+\cdots+y_{n-1}
\]
代入 $fu=gv$ 并比较系数,得到一个 $m+n$ 元齐次线性方程组.
其系数矩阵的行列式即引出结式.

\begin{definition}
上述两个多项式的结式(Sylvester 行列式)记作
\[
  R(f,g)=\det S(f,g),
\]
其中 Sylvester 矩阵为
\[
S(f,g)=
\begin{pmatrix}
 a_0&a_1&\cdots&a_n&0&\cdots&0\\
 0&a_0&a_1&\cdots&a_n&\ddots&\vdots\\
 \vdots&\ddots&\ddots&\ddots&&\ddots&0\\
 0&\cdots&0&a_0&a_1&\cdots&a_n\\
 b_0&b_1&\cdots&b_m&0&\cdots&0\\
 0&b_0&b_1&\cdots&b_m&\ddots&\vdots\\
 \vdots&\ddots&\ddots&\ddots&&\ddots&0\\
 0&\cdots&0&b_0&b_1&\cdots&b_m
\end{pmatrix},
\]
其中前 $m$ 行由 $f$ 的系数错位排列,后 $n$ 行由 $g$ 的系数错位排列.
\end{definition}

\begin{theorem}[\markstar]
\[
  (f,g)=1\iff R(f,g)\ne0,
\]
等价地,
\[
  (f,g)\ne1\iff R(f,g)=0.
\]
\end{theorem}

\begin{theorem}[\markstar]
若 $f$ 的根为 $x_1,\ldots,x_n$,$g$ 的根为 $y_1,\ldots,y_m$,则
\[
\begin{aligned}
  R(f,g)
  &=a_0^m b_0^n
    \prod\limits_{j=1}^{n}\prod\limits_{i=1}^{m}(x_j-y_i)\\
  &=a_0^m\prod\limits_{j=1}^{n}g(x_j)\\
  &=(-1)^{mn}b_0^n\prod\limits_{i=1}^{m}f(y_i).
\end{aligned}
\]
\end{theorem}

\begin{definition}
$n$ 次多项式
\[
  f(x)=a_0x^n+a_1x^{n-1}+\cdots+a_n
\]
的判别式定义为
\[
  \Delta(f)
  =(-1)^{\frac{n(n-1)}2}a_0^{-1}R(f,f').
\]
\end{definition}

\begin{theorem}[\markstar]
若 $x_1,\ldots,x_n$ 为 $f$ 的全部根,则
\[
  \Delta(f)
  =a_0^{2n-2}
  \prod\limits_{1\le i<j\le n}(x_i-x_j)^2.
\]
因此
\[
  f(x)\text{ 有重根}
  \iff
  \Delta(f)=0.
\]
\end{theorem}

%例题1.6.1
\begin{example}
求下列结式:
\begin{enumerate}
  \item $x^n+x+1$ 与 $x^2-3x+2$;
  \item $x^n+1$ 与 $(x-1)^n$;
  \item $x^n-1$ 与 $x^m-1$.
\end{enumerate}
\end{example}
\exampleblank

\begin{solution}

(1) 因
\[
  x^2-3x+2=(x-1)(x-2),
\]
由结式的根表达式
\[
  R(x^n+x+1,x^2-3x+2)
  =f(1)f(2)
  =3(2^n+3).
\]

(2) 第二个多项式只有根 $1$,重数为 $n$.因此
\[
  R(x^n+1,(x-1)^n)
  =(-1)^{n^2}[1^n+1]^n
  =(-1)^n2^n.
\]

(3) 两个多项式都有公共根 $1$,故
\[
  R(x^n-1,x^m-1)=0.
\]
\end{solution}

%例题1.6.2
\begin{example}
求判别式:
\begin{enumerate}
  \item $x^n+ax+b$;
  \item $x^n+ax^k+b$.
\end{enumerate}
\end{example}
\exampleblank

\begin{solution}
由
\[
  \Delta(f)=(-1)^{\frac{n(n-1)}2}R(f,f')
\]
及结式乘法性质可得:

(1)
\[
\boxed{
  \Delta(x^n+ax+b)
  =(-1)^{\frac{n(n-1)}2}
  \left[n^n b^{n-1}
  +(-1)^{n-1}(n-1)^{n-1}a^n\right].}
\]

(2) 令
\[
  d=(n,k).
\]
则一般情形为
\[
\boxed{
\begin{aligned}
  \Delta(x^n+ax^k+b)
  ={}&(-1)^{\frac{n(n-1)}2}b^{k-1}\\
  &\times\left[
    n^{\frac{n}{d}}b^{\frac{n-k}{d}}
    -(-1)^{\frac{n}{d}}
    (n-k)^{\frac{n-k}{d}}k^{\frac{k}{d}}a^{\frac{n}{d}}
  \right]^d.
\end{aligned}}
\]
特别地,当 $(n,k)=1$ 时,上式化为
\[
  \Delta
  =(-1)^{\frac{n(n-1)}2}b^{k-1}
  \left[n^n b^{n-k}
  +(-1)^{n-1}(n-k)^{n-k}k^k a^n\right].
\]
\end{solution}

%例题1.6.3
\begin{example}
解方程组:
\begin{enumerate}
  \item
  \[
  \begin{cases}
    x^3+y^3=7(x+y),\\
    x^2+y^2=13;
  \end{cases}
  \]
  \item
  \[
  \begin{cases}
    -ay+x(1-x^2-y^2)=0,\\
    ax+y(1-x^2-y^2)=a.
  \end{cases}
  \]
\end{enumerate}
\end{example}
\exampleblank

\begin{solution}

(1) 令
\[
  s=x+y,
  \qquad
  p=xy.
\]
由 $x^2+y^2=13$ 得
\[
  p=\dfrac{s^2-13}{2}.
\]
又
\[
  x^3+y^3=s^3-3ps=7s,
\]
代入得到
\[
  s(25-s^2)=0.
\]
故 $s=0,\pm5$.逐一求根,得到
\[
\boxed{
\begin{gathered}
  (x,y)=(2,3),(3,2),(-2,-3),(-3,-2),\\
  \left(\dfrac{\sqrt{26}}2,-\dfrac{\sqrt{26}}2\right),
  \left(-\dfrac{\sqrt{26}}2,\dfrac{\sqrt{26}}2\right).
\end{gathered}}
\]

(2) 设
\[
  \rho=x^2+y^2,
  \qquad
  z=x+iy.
\]
两式可合并为
\[
  (1-\rho+ia)z=ia.
\]
若 $a\ne0$,取模得到
\[
  \rho\left((1-\rho)^2+a^2\right)=a^2,
\]
即
\[
  (\rho-1)(\rho^2-\rho+a^2)=0.
\]
因此恒有解
\[
  (x,y)=(1,0).
\]
若 $|a|\le\dfrac12$,还可取
\[
  \rho_{\pm}=\dfrac{1\pm\sqrt{1-4a^2}}2,
\]
并由原方程得
\[
  (x,y)=(\rho_{\pm},a).
\]
故当 $a\ne0$ 时:
\[
\boxed{
\begin{cases}
|a|>\dfrac12:&(x,y)=(1,0),\\[0.6ex]
0<|a|\le\dfrac12:&(x,y)=(1,0),\ (\rho_+,a),\ (\rho_-,a),
\end{cases}}
\]
其中 $|a|=\dfrac12$ 时 $\rho_+=\rho_-$,重复解只计一次.

若 $a=0$,原方程化为
\[
  x(1-x^2-y^2)=0,
  \qquad
  y(1-x^2-y^2)=0,
\]
故解为
\[
  (0,0)
  \quad\text{以及}\quad
  x^2+y^2=1
\]
上的所有实点.
\end{solution}

%例题1.6.4
\begin{example}
判定以下方程组是否有有理解:
\[
\begin{cases}
  y^2+2x^2y-1=0,\\
  6x^2-y^2-3y=0.
\end{cases}
\]
\end{example}
\exampleblank

\begin{solution}
由第二式
\[
  x^2=\dfrac{y^2+3y}{6}.
\]
代入第一式得
\[
  y^3+6y^2-3=0.
\]
若 $x,y\in\mathbb{Q}$,则 $y$ 必是上述首一整系数多项式的有理根,因此只能是整数且整除 $3$,
即候选为 $\pm1,\pm3$.逐一代入均不为零,矛盾.因此该方程组无有理解.
\end{solution}

%例题1.6.5
\begin{example}
求参数曲线
\[
\begin{cases}
  x=\dfrac{2(t+1)}{t^2+1},\\[1ex]
  y=\dfrac{t^2}{2t-1}
\end{cases}
\]
的直角坐标方程.
\end{example}
\exampleblank

\begin{solution}
由
\[
  x(t^2+1)-2(t+1)=0,
\]
\[
  y(2t-1)-t^2=0
\]
消去参数 $t$,取两多项式关于 $t$ 的结式,得到
\[
\boxed{
  5x^2y^2-2x^2y+x^2-12xy^2-4x+12y+4=0.}
\]
参数表达式还要求 $t\ne\dfrac12$;上式给出参数曲线所落在的代数曲线.
\end{solution}
